\documentclass{article}
\usepackage[margin=1in]{geometry}
\usepackage{amsmath,amssymb,amsfonts}
\usepackage{graphicx,algorithm,algpseudocode,booktabs,listings}
\usepackage{cite}

\title{\bf Commit-Time Governance: Enforcing Recoverability at the Point of Action}
\author{[Authors anonymized for submission]}

\begin{document}
\maketitle

%%%%%%%%%%%%%%%%%%%%%%%%%%%%%%%%%%%%%%%%%%%%%%%%%%%%%%%%%%%%%%%%%%%%%%%%%%%%%%%%
\begin{abstract}
We introduce commit-time governance: a minimal enforcement mechanism that checks whether an action can be recovered from at the exact moment it becomes irreversible. Current systems separate governance into pre-action policy and post-action audit, leaving the commit — the point of no return — ungoverned. We close this gap with four contributions: (1) a commit admissibility condition based on recoverability, not just static legitimacy; (2) an executable gate with reproducible traces demonstrating enforcement behavior; (3) a barrier-certificate proof that the gate is sound with respect to a certified safe subset; and (4) an adversarial robustness analysis showing fail-closed degradation under bounded attacks. The framework operates on a 4-dimensional state space (governance, control, autonomy, trust), executes in constant time, and enforces recoverability strictly at the point of action.
\end{abstract}

%%%%%%%%%%%%%%%%%%%%%%%%%%%%%%%%%%%%%%%%%%%%%%%%%%%%%%%%%%%%%%%%%%%%%%%%%%%%%%%%
\section{Introduction}

Governance systems typically enforce before action (policy, guardrails) or after action (monitoring, audit). Both miss the critical moment: the commit, when intention becomes irreversible effect.

We define \textbf{commit-time governance}: an action is allowed if and only if, after it executes, the system remains in a state from which recovery is possible. This shifts governance from static legitimacy to dynamic recoverability.

Our contributions are minimal and sequential:
\begin{enumerate}
    \item \textbf{Commit admissibility}: $\Phi(\mathbf{x},u)\in\mathcal{R}{rob}(\tau)$
    \item \textbf{Executable gate}: constant-time implementation with traces
    \item \textbf{Certified approximation}: barrier-certificate soundness
    \item \textbf{Adversarial robustness}: fail-closed degradation
\end{enumerate}

\textbf{Scope.} The framework governs individual commits. It does not guarantee optimal trajectories, global safety, or trajectory-level completeness.

%%%%%%%%%%%%%%%%%%%%%%%%%%%%%%%%%%%%%%%%%%%%%%%%%%%%%%%%%%%%%%%%%%%%%%%%%%%%%%%%
\section{Minimal Formal Setup}

\textbf{State.} $\mathbf{x}=(g,c,a,t)\in[0,1]^4$.

\textbf{Legitimacy.}
[
\mathcal{A}={\mathbf{x}:m(\mathbf{x})\geq 0}, \quad m(\mathbf{x})=\Lambda(\mathbf{x})-a
]
[
\Lambda(\mathbf{x})=Kg^\alpha c^\beta t^\gamma
]

\textbf{Recoverability.}
[
\mathcal{R}(0)={\text{states recoverable to }\mathcal{A}}, \quad
\mathcal{R}{rob}(\tau)={\text{states recoverable under lag }\tau}
]

\textbf{Normalized geometry.}
[
\tilde{\mathbf{x}}=\frac{\mathbf{x}}{|\mathbf{x}|1}, \quad \mathbf{R}=(1/4,1/4,1/4,1/4)
]
[
d_F=\min_i \tilde{x}i, \quad d\Delta=|\tilde{\mathbf{x}}-\mathbf{R}|2
]

\textbf{Lag.} $\tau=\tau{obs}+\tau{dec}+\tau_{act}$.

%%%%%%%%%%%%%%%%%%%%%%%%%%%%%%%%%%%%%%%%%%%%%%%%%%%%%%%%%%%%%%%%%%%%%%%%%%%%%%%%
\section{The Commit Gate}

\textbf{Commit operator.}
[
\Phi(\mathbf{x},u)=\text{clip}{[0,1]^4}(\mathbf{x}+u)
]

\textbf{Ideal admissibility.}
[
\text{ALLOW}{ideal}(u;\mathbf{x},\tau)\Leftrightarrow\Phi(\mathbf{x},u)\in\mathcal{R}{rob}(\tau)
]

\textbf{Commit Safety Theorem.}
If all executed actions satisfy $\text{ALLOW}{ideal}$, the system never enters a state from which collapse is unavoidable under bounded disturbances.

\textit{Proof sketch.} Induction on commit events preserves recoverability. \hfill $\blacksquare$

\textbf{Practical constraint.}
Exact evaluation of $\mathcal{R}{rob}(\tau)$ requires online reachability analysis.

\textbf{Approach.}
Use a conservative certified subset.

%%%%%%%%%%%%%%%%%%%%%%%%%%%%%%%%%%%%%%%%%%%%%%%%%%%%%%%%%%%%%%%%%%%%%%%%%%%%%%%%
\section{Executable Evidence}

We implement: state, commit operator, certified test, and ALLOW/DENY gate.

\textbf{Certified subset.}
[
\mathcal{S}{cert}(\tau)={\mathbf{x}:B_F>0 \wedge B_D>0}
]

[
B_F=d_F-\epsilon_{crit}(\tau), \quad
B_D=d_{max}(\tau)-d_\Delta
]

\textbf{Guarantee.}
Membership in $\mathcal{S}{cert}(\tau)$ is sufficient (not necessary) for membership in $\mathcal{R}{rob}(\tau)$.

\begin{table}[h]
\centering
\caption{Gate Performance}
\begin{tabular}{lcccc}
Scenario & Actions & Allowed & Denied & Recovery \
\hline
Drone & 4 & 3 & 1 & 100% \
Human & 9 & 8 & 1 & 100% \
Consensus & 11 & 10 & 1 & 100% \
\end{tabular}
\end{table}

\textbf{Key finding.}
In each scenario, the gate denies the transition that violates the certified safety conditions. No allowed transition produces a state outside $\mathcal{S}{cert}(\tau)$.

%%%%%%%%%%%%%%%%%%%%%%%%%%%%%%%%%%%%%%%%%%%%%%%%%%%%%%%%%%%%%%%%%%%%%%%%%%%%%%%%
\section{Certified Approximation}

\textbf{Calibration assumptions.}
\begin{itemize}
    \item C1: $d_F>\epsilon{crit}$ implies lag-safe separation
    \item C2: $d_\Delta<d_{max}$ implies recoverability to $\mathcal{A}$
\end{itemize}

\textbf{Soundness.}
[
\mathcal{S}{cert}(\tau)\subseteq\mathcal{R}{rob}(\tau)
]

\textbf{Corollary.}
[
\text{ALLOW}{impl}(u;\mathbf{x},\tau)\Rightarrow\text{ALLOW}{ideal}(u;\mathbf{x},\tau)
]

\textbf{Approximation error.}
[
\mathcal{E}(\tau)=\mathcal{R}{rob}(\tau)\setminus\mathcal{S}{cert}(\tau)
]

%%%%%%%%%%%%%%%%%%%%%%%%%%%%%%%%%%%%%%%%%%%%%%%%%%%%%%%%%%%%%%%%%%%%%%%%%%%%%%%%
\section{Adversarial Robustness}

\textbf{Threats.}
State perturbation, lag inflation, boundary surfing, recovery disruption.

\textbf{Bounded attacks.}
\begin{itemize}
    \item State deception: expands certified subset
    \item Lag inflation: stricter evaluation (fail-closed)
    \item Boundary surfing: pointwise safe
    \item Recovery disruption: bounded stability preserved
\end{itemize}

\textbf{Unbounded attacks.}
System degrades to denial or safe mode.

\textbf{Property.}
Fail-closed behavior: failure results in denial, not unsafe execution.

%%%%%%%%%%%%%%%%%%%%%%%%%%%%%%%%%%%%%%%%%%%%%%%%%%%%%%%%%%%%%%%%%%%%%%%%%%%%%%%%
\section{Discussion}

\textbf{Guarantees.}
\begin{itemize}
    \item Allowed commits remain recoverable
    \item Unsafe commits are denied pre-execution
    \item Soundness holds under bounded perturbations
\end{itemize}

\textbf{Limitations.}
\begin{itemize}
    \item Not trajectory-complete
    \item Not globally optimal
    \item Requires external monitoring for sequence effects
\end{itemize}

%%%%%%%%%%%%%%%%%%%%%%%%%%%%%%%%%%%%%%%%%%%%%%%%%%%%%%%%%%%%%%%%%%%%%%%%%%%%%%%%
\section{Conclusion}

Commit-time governance enforces recoverability at the moment of action. It is a minimal, constant-time mechanism that ensures no allowed transition leads to unavoidable collapse. The framework is intentionally narrow: it governs commits, not trajectories.

\end{document}
